\section{Discussion}
\label{sec:discussion}

\subsection{Interpreting the Ceiling Effect}

The central finding of this work is a null result: evolved prompts do not outperform a carefully crafted human baseline.
We argue this is informative rather than negative.
It establishes that for frontier models on simple tasks, goal-persistence \emph{can} be instructed into existence---the hypothesis that ``desire must be evolved rather than instructed'' is not supported at this level of task difficulty.
This finding has practical implications: for current deployments of frontier models on well-understood tasks, investing in prompt engineering is sufficient and the overhead of evolutionary search is unnecessary.

The ceiling effect likely reflects two factors.
First, \gptmini is a capable model that responds well to explicit instructions.
\citet{arike2025goaldrift} similarly found near-zero drift for Claude~3.5 Sonnet with strong elicitation prompts, suggesting this is a property of frontier models generally.
Second, the counting task is simple enough that any sufficiently explicit prompt can prevent drift.
The relevant question is not whether evolution helps here, but at what level of task complexity evolution \emph{begins} to outperform human design.

\subsection{What Evolution Teaches Us}

Despite the performance tie, the qualitative analysis reveals three insights about what makes prompts goal-persistent.

{\bf Structure matters.}
The convergence of all evolved prompts to GOAL/RULE/PRIORITY/MANDATE organization suggests that structured formatting is a general-purpose persistence strategy.
This is supported by the fact that a single seed prompt with this structure achieved perfect fitness even before evolution began (Prompt~4, generation~0).

{\bf Recovery is distinct from prevention.}
Evolved prompts independently discover recovery language (``if interrupted, resume''), treating drift recovery as a separate concern from drift prevention.
Human prompt designers, including the author of the \strongelicit baseline, focus on preventing drift but do not explicitly instruct recovery.
This suggests that for harder tasks where drift occasionally succeeds, evolved prompts may provide an advantage through their recovery mechanisms.

{\bf Redundancy provides robustness.}
The 4$\times$ expansion in prompt length is not noise---it reflects a strategy of expressing the same instruction in multiple complementary forms.
This redundancy may act as a form of implicit self-reinforcement, analogous to how repeated priming in psychology strengthens goal activation \citep{bargh2001automaticity}.

\subsection{Limitations}

\para{Ceiling effect prevents primary hypothesis test.}
The most significant limitation is that the ceiling effect prevents us from testing whether evolution produces meaningfully better prompts.
Our conclusion---that evolution matches but does not exceed human design---is contingent on this specific task/model combination.

\para{Single model.}
We test only on \gptmini.
Weaker models (\eg GPT-3.5, open-weight models with 7--8B parameters) may show larger gaps between evolved and human prompts, as they are more sensitive to prompt formulation \citep{arike2025goaldrift}.

\para{Simple task.}
Counting is trivially verifiable but may not represent real-world goal-persistence challenges such as maintaining a complex persona, following a multi-step plan, or managing tool use across long horizons.

\para{Fixed distractions.}
Evolution optimizes against a fixed set of 5 training distractions.
A co-evolutionary approach \citep{aegis2025} that evolves distractors alongside prompts would produce more robust results and avoid ceiling effects by escalating difficulty.

\para{LLM-based mutation bias.}
Using \gptmini for both mutation and evaluation may bias evolution toward prompts the model finds natural to generate, potentially limiting the discovery of unconventional but effective strategies.

\para{Conversation length.}
All scenarios use 20 turns.
Real goal drift often emerges only in longer conversations (50--100+ turns) where the initial instruction's influence decays.

\subsection{Broader Implications}

\para{For LLM deployment.}
Our results suggest a practical hierarchy for prompt engineering: invest in structured, emphatic prompts with explicit distraction negation before considering evolutionary optimization.
The feature analysis (\tabref{tab:features}) provides a concrete checklist of properties that goal-persistent prompts should contain.

\para{For AI alignment.}
The question of whether goal-directedness must be evolved or can be instructed connects to fundamental questions in AI alignment about the nature of agency and goals.
Our results suggest that for current systems, instruction is sufficient---but this may change as models become more autonomous and operate over longer horizons.

\para{For prompt evolution research.}
The rapid convergence (2 generations) and population homogenization we observe suggest that DE-based evolution may be over-parameterized for simple fitness landscapes.
Adaptive population sizing or early stopping could reduce computational cost without sacrificing solution quality.
